% --- CORRECCION: Forzar interpretacion raw de la entrada ---
\UseRawInputEncoding
\documentclass[aspectratio=169, 10pt, xcolor=table]{beamer}

% --- Configuracion del Tema ---
\usetheme{Madrid}
\usecolortheme{dolphin}

% --- Paquetes necesarios ---
\usepackage[utf8]{inputenc}
\usepackage[T1]{fontenc}
\usepackage[spanish,es-noquoting]{babel}
\usepackage{amsmath, amssymb}
\usepackage{mathtools}
\usepackage{booktabs}
\usepackage{graphicx}
\usepackage{listings}
\usepackage{tikz}
\usetikzlibrary{arrows.meta, positioning, shapes.geometric, babel}
\usepackage{etoolbox}
% \usepackage{upquote} % Eliminado

% --- Estilo para codigo SQL, bash, Python, YAML ---
\definecolor{codegreen}{rgb}{0,0.6,0}
\definecolor{codegray}{rgb}{0.5,0.5,0.5}
\definecolor{codepurple}{rgb}{0.58,0,0.82}
\definecolor{backcolour}{rgb}{0.95,0.95,0.92}
\lstdefinestyle{bashstyle}{
    language=bash,
    backgroundcolor=\color{backcolour},
    commentstyle=\color{codegreen},
    keywordstyle=\color{blue},
    numberstyle=\tiny\color{codegray},
    stringstyle=\color{codepurple},
    basicstyle=\ttfamily\scriptsize,
    breaklines=true,
    frame=single,
    showstringspaces=false,
    literate={á}{{\'a}}1 {é}{{\'e}}1 {í}{{\'i}}1 {ó}{{\'o}}1 {ú}{{\'u}}1 {ñ}{{\~n}}1 {ü}{{\"u}}1
}
\lstdefinestyle{sqlstyle}{
    language=SQL,
    backgroundcolor=\color{backcolour},
    commentstyle=\color{codegreen},
    keywordstyle=\color{blue},
    numberstyle=\tiny\color{codegray},
    stringstyle=\color{codepurple},
    basicstyle=\ttfamily\scriptsize,
    breaklines=true,
    frame=single,
    showstringspaces=false,
    literate={á}{{\'a}}1 {é}{{\'e}}1 {í}{{\'i}}1 {ó}{{\'o}}1 {ú}{{\'u}}1 {ñ}{{\~n}}1 {ü}{{\"u}}1
}
\lstdefinestyle{pythonstyle}{
    language=Python,
    backgroundcolor=\color{backcolour},
    commentstyle=\color{codegreen},
    keywordstyle=\color{blue},
    numberstyle=\tiny\color{codegray},
    stringstyle=\color{codepurple},
    basicstyle=\ttfamily\scriptsize,
    breaklines=true,
    frame=single,
    showstringspaces=false,
    literate={á}{{\'a}}1 {é}{{\'e}}1 {í}{{\'i}}1 {ó}{{\'o}}1 {ú}{{\'u}}1 {ñ}{{\~n}}1 {ü}{{\"u}}1
}
\lstdefinestyle{yamlstyle}{
    language=bash, % no hay lenguaje YAML nativo, se usa bash como base
    backgroundcolor=\color{backcolour},
    commentstyle=\color{codegreen},
    keywordstyle=\color{blue},
    numberstyle=\tiny\color{codegray},
    stringstyle=\color{codepurple},
    basicstyle=\ttfamily\scriptsize,
    breaklines=true,
    frame=single,
    showstringspaces=false,
    literate={á}{{\'a}}1 {é}{{\'e}}1 {í}{{\'i}}1 {ó}{{\'o}}1 {ú}{{\'u}}1 {ñ}{{\~n}}1 {ü}{{\"u}}1
}
\lstset{style=sqlstyle} % Por defecto SQL

% --- Pie de pagina personalizado ---
\makeatletter
\setbeamertemplate{footline}{
  \leavevmode%
  \hbox{%
  \begin{beamercolorbox}[wd=.333333\paperwidth,ht=2.25ex,dp=1ex,center]{author in head/foot}%
    \usebeamerfont{author in head/foot}\insertshortauthor
  \end{beamercolorbox}%
  \begin{beamercolorbox}[wd=.333333\paperwidth,ht=2.25ex,dp=1ex,center]{title in head/foot}%
    \usebeamerfont{title in head/foot}Sesión 14
  \end{beamercolorbox}%
  \begin{beamercolorbox}[wd=.333333\paperwidth,ht=2.25ex,dp=1ex,right]{date in head/foot}%
    \usebeamerfont{date in head/foot}BBDD \hfill \insertframenumber/\inserttotalframenumber \hspace*{0.3cm}
  \end{beamercolorbox}}%
  \vskip0pt%
}
\makeatother

% --- Datos de la portada ---
\title[SESIÓN 14 BBDD]{\textbf{Sesión 14}}
\subtitle{Gestión de Errores, Recuperación y Arquitecturas Híbridas}
\author[Prof. C. Sánchez]{Carlos César Sánchez Coronel}
\institute{\textbf{BASES DE DATOS}}
\date{2026}

\setbeamertemplate{navigation symbols}{}

\AtBeginSection[]{
  \begin{frame}
    \frametitle{Contenido de la sección}
    \tableofcontents[currentsection]
  \end{frame}
}

\begin{document}

% --- Ya no usamos \shorthandoff{<>} porque con es-noquoting no son shorthands activos ---

\begin{frame}
    \titlepage
\end{frame}

\begin{frame}{Contenido general}
    \tableofcontents
\end{frame}

% ============================================================================
% SECCIÓN 1: INTRODUCCIÓN
% ============================================================================
\section{Introducción}

\begin{frame}{Objetivos de la sesión}
    \begin{itemize}
        \item Conocer las operaciones de alto riesgo y las mejores prácticas para evitarlas.
        \item Entender los mecanismos de recuperación ante desastres: backups, snapshots, point-in-time recovery.
        \item Analizar los desafíos específicos de la nube: latencia, costos de egreso, data drift.
        \item Diseñar arquitecturas híbridas que integren motores SQL, NoSQL y vectoriales.
        \item Aprender a orquestar múltiples bases de datos con contenedores e infraestructura como código.
        \item Implementar sincronización entre sistemas mediante CDC y triggers.
        \item Resolver ejercicios prácticos que integran todos los conceptos.
    \end{itemize}
\end{frame}

% ============================================================================
% SECCIÓN 2: OPERACIONES DE ALTO RIESGO Y PREVENCIÓN
% ============================================================================
\section{Operaciones de Alto Riesgo y Prevención}

\begin{frame}{DROP vs TRUNCATE vs DELETE}
    \begin{table}
        \centering
        \begin{tabular}{l|c|c}
            \toprule
            Operación & ¿Elimina estructura? & ¿Registra por fila? \\
            \midrule
            DROP & Sí & No (metadata) \\
            TRUNCATE & No & No (desasigna páginas) \\
            DELETE sin WHERE & No & Sí (puede ser lento) \\
            \bottomrule
        \end{tabular}
    \end{table}
    \begin{itemize}
        \item DROP: irreversible sin backup.
        \item TRUNCATE: rápido, pero no se puede deshacer con ROLLBACK en muchos motores.
        \item DELETE: transaccional, pero genera muchos logs.
    \end{itemize}
\end{frame}

\begin{frame}{Buenas prácticas para prevenir desastres}
    \begin{itemize}
        \item Siempre realizar un \texttt{SELECT} antes de un \texttt{DELETE} o \texttt{UPDATE} para verificar la condición.
        \item Usar transacciones: \texttt{BEGIN; DELETE ...; ROLLBACK;} para probar.
        \item En entornos críticos, revocar permisos de \texttt{DROP} y \texttt{TRUNCATE} para usuarios regulares.
        \item Implementar \textbf{soft delete} (columna \texttt{activo}) en lugar de eliminación física.
        \item Usar \texttt{sql\_safe\_updates} en MySQL (requiere WHERE con key).
        \item Auditoría de consultas con herramientas como \texttt{pgaudit}.
    \end{itemize}
\end{frame}

\begin{frame}[fragile]{Ejemplo: Soft Delete}
    \begin{lstlisting}[style=sqlstyle]
ALTER TABLE productos ADD COLUMN activo BOOLEAN DEFAULT TRUE;

-- En lugar de DELETE
UPDATE productos SET activo = FALSE WHERE id = 123;

-- Las consultas ignoran productos inactivos
SELECT * FROM productos WHERE activo = TRUE;
    \end{lstlisting}
\end{frame}

% ============================================================================
% SECCIÓN 3: MECANISMOS DE RESILIENCIA Y RECUPERACIÓN ANTE DESASTRES
% ============================================================================
\section{Recuperación ante Desastres (DRP)}

\begin{frame}{RTO y RPO}
    \begin{description}
        \item[RTO] Recovery Time Objective: tiempo máximo para restaurar el servicio.
        \item[RPO] Recovery Point Objective: pérdida máxima de datos aceptable (expresada en tiempo).
    \end{description}
    \begin{itemize}
        \item RPO bajo requiere replicación continua o backups incrementales frecuentes.
        \item RTO bajo requiere procedimientos automatizados (failover automático).
    \end{itemize}
\end{frame}

\begin{frame}{Backups y Snapshots}
    \begin{description}
        \item[Backups lógicos] Volcados de datos en SQL (ej. \texttt{pg\_dump}). Portables, lentos.
        \item[Backups físicos] Copias del sistema de archivos (snapshots EBS). Rápidos, dependientes del proveedor.
        \item[Backups incrementales] Solo guardan cambios desde el último backup. Ahorran espacio.
    \end{description}
\end{frame}

\begin{frame}{Point-in-Time Recovery (PITR)}
    \begin{itemize}
        \item Permite restaurar la base a un momento exacto antes de un incidente.
        \item Requiere:
        \begin{itemize}
            \item Backup completo (basebackup).
            \item Archivos WAL (Write-Ahead Logs) desde ese backup.
            \item Configuración de archiving y recovery target.
        \end{itemize}
    \end{itemize}
\end{frame}

\begin{frame}[fragile]{Ejemplo: Configuración de PITR en PostgreSQL}
    \begin{lstlisting}[style=bashstyle]
# Configurar archiving en postgresql.conf
wal_level = replica
archive_mode = on
archive_command = 'cp %p /wal_archive/%f'
archive_timeout = 60

# Backup base
pg_basebackup -D /backup/base -F t -z -P

# En postgresql.auto.conf para recuperación:
restore_command = 'cp /wal_archive/%f %p'
recovery_target_time = '2025-03-20 14:30:00'
    \end{lstlisting}
    Al iniciar PostgreSQL, aplica WAL hasta ese momento.
\end{frame}

% ============================================================================
% SECCIÓN 4: DESAFÍOS EN LA NUBE Y DATA DRIFT
% ============================================================================
\section{Desafíos en la Nube y Data Drift}

\begin{frame}{Latencia y Costos de Egreso}
    \begin{itemize}
        \item Mover datos entre regiones o fuera del proveedor tiene costo (egress).
        \item Estrategias:
        \begin{itemize}
            \item Mantener datos en la misma región.
            \item Comprimir antes de transferir.
            \item Usar redes privadas (VPC peering).
        \end{itemize}
    \end{itemize}
\end{frame}

\begin{frame}{Data Drift (Deriva de Datos)}
    \begin{itemize}
        \item Cambio en la distribución de los datos de entrada que degrada modelos de IA.
        \item Se detecta comparando estadísticas (media, desviación) o tests como Kolmogorov-Smirnov.
        \item Herramientas: Evidently, Great Expectations, WhyLabs.
    \end{itemize}
\end{frame}

\begin{frame}[fragile]{Detección de data drift con Python}
    \begin{lstlisting}[style=pythonstyle]
import numpy as np
from scipy.stats import ks_2samp

train = np.random.normal(100, 15, 1000)
prod = np.random.normal(105, 18, 500)

stat, p_value = ks_2samp(train, prod)
alpha = 0.05
if p_value < alpha:
    print("Drift detectado (p=%.4f)" % p_value)
else:
    print("No hay evidencia de drift")
    \end{lstlisting}
\end{frame}

% ============================================================================
% SECCIÓN 5: ARQUITECTURAS HÍBRIDAS
% ============================================================================
\section{Arquitecturas Híbridas}

\begin{frame}{Capas en una arquitectura políglota}
    \begin{description}
        \item[Transaccional (SQL)] Datos maestros, ACID. Ej: PostgreSQL.
        \item[Velocidad (NoSQL)] Caché, sesiones, baja latencia. Ej: Redis.
        \item[Semántica (Vectorial)] Búsqueda por similitud, embeddings. Ej: pgvector, Milvus.
    \end{description}
\end{frame}

\begin{frame}[fragile]{Orquestación con Docker Compose}
    \begin{lstlisting}[style=yamlstyle]
version: '3'
services:
  postgres:
    image: postgres:15
    environment:
      POSTGRES_PASSWORD: admin
      POSTGRES_DB: test
    volumes:
      - pgdata:/var/lib/postgresql/data
    networks:
      - app-net
  redis:
    image: redis:7
    networks:
      - app-net
  app:
    build: ./app
    depends_on:
      - postgres
      - redis
    networks:
      - app-net
networks:
  app-net:
    driver: bridge
volumes:
  pgdata:
    \end{lstlisting}
\end{frame}

\begin{frame}{Infraestructura como Código (IaC)}
    \begin{itemize}
        \item Definir infraestructura de forma declarativa (Terraform, CloudFormation).
        \item Ventajas: replicable, versionable, evita errores manuales.
    \end{itemize}
    \begin{exampleblock}{Ejemplo Terraform para RDS}
        \begin{lstlisting}[style=bashstyle]

        \end{lstlisting}
    \end{exampleblock}
\end{frame}

% ============================================================================
% SECCIÓN 6: SINCRONIZACIÓN ENTRE MOTORES
% ============================================================================
\section{Sincronización y Consistencia}

\begin{frame}[fragile]{Sincronización con Triggers en SQL}
    \begin{lstlisting}[style=sqlstyle]
CREATE TABLE cambios_productos (
    id SERIAL PRIMARY KEY,
    producto_id INT,
    nombre TEXT,
    operacion CHAR(1),
    procesado BOOLEAN DEFAULT FALSE
);

CREATE OR REPLACE FUNCTION log_producto() RETURNS TRIGGER AS $$
BEGIN
    INSERT INTO cambios_productos (producto_id, nombre, operacion)
    VALUES (NEW.id, NEW.nombre, TG_OP);
    RETURN NEW;
END;
$$ LANGUAGE plpgsql;

CREATE TRIGGER trig_productos
AFTER INSERT OR UPDATE OR DELETE ON productos
FOR EACH ROW EXECUTE FUNCTION log_producto();
    \end{lstlisting}
    Un worker externo lee la tabla y actualiza Redis/vectorial.
\end{frame}

\begin{frame}[fragile]{Worker Python para sincronizar con Redis}
    \begin{lstlisting}[style=pythonstyle]
import psycopg2, redis, time

conn_pg = psycopg2.connect(dbname="test", user="postgres")
cur_pg = conn_pg.cursor()
r = redis.Redis(host='localhost', port=6379, decode_responses=True)

while True:
    cur_pg.execute("SELECT id, producto_id, nombre, operacion FROM cambios_productos WHERE procesado = FALSE LIMIT 1")
    row = cur_pg.fetchone()
    if row:
        cambio_id, prod_id, nombre, op = row
        if op in ('INSERT', 'UPDATE'):
            r.set(f"producto:{prod_id}", nombre)
        elif op == 'DELETE':
            r.delete(f"producto:{prod_id}")
        cur_pg.execute("UPDATE cambios_productos SET procesado = TRUE WHERE id = %s", (cambio_id,))
        conn_pg.commit()
        print(f"Sincronizado {prod_id}")
    time.sleep(1)
    \end{lstlisting}
\end{frame}

\begin{frame}{Change Data Capture (CDC)}
    \begin{center}
        \begin{tikzpicture}[node distance=2cm, auto]
            \node (pg) [draw] {PostgreSQL};
            \node (debezium) [draw, right of=pg] {Debezium};
            \node (kafka) [draw, right of=debezium] {Kafka};
            \node (consumer) [draw, right of=kafka] {Consumidor};
            \node (milvus) [draw, right of=consumer] {Milvus};
            \draw[->] (pg) -- (debezium) node[midway, above] {log};
            \draw[->] (debezium) -- (kafka) node[midway, above] {eventos};
            \draw[->] (kafka) -- (consumer);
            \draw[->] (consumer) -- (milvus);
        \end{tikzpicture}
    \end{center}
    \begin{itemize}
        \item Debezium captura cambios del log de transacciones.
        \item Publica en Kafka; consumidores actualizan bases vectoriales.
    \end{itemize}
\end{frame}

\begin{frame}{Consistencia Eventual}
    \begin{itemize}
        \item La replicación asíncrona introduce latencia.
        \item Aplicaciones deben tolerar que la búsqueda no incluya los últimos cambios.
        \item Estrategias: marcar datos como "no indexados" hasta sincronizar, usar TTL.
    \end{itemize}
\end{frame}

% ============================================================================
% SECCIÓN 7: EJERCICIOS RESUELTOS
% ============================================================================
\section{Ejercicios Resueltos}

\begin{frame}[fragile]{Ejercicio 1: Prevención de DROP accidental}
    \textbf{Enunciado:} En PostgreSQL, crear un trigger que impida ejecutar DROP TABLE en una tabla crítica.
    \begin{block}{Solución con Event Trigger}
        \begin{lstlisting}[style=sqlstyle]
CREATE OR REPLACE FUNCTION abort_drop()
RETURNS event_trigger AS $$
BEGIN
    RAISE EXCEPTION 'No se permite DROP en esta base de datos';
END;
$$ LANGUAGE plpgsql;

CREATE EVENT TRIGGER abort_drop_trigger
ON sql_drop
EXECUTE FUNCTION abort_drop();
        \end{lstlisting}
        (Nota: impide cualquier DROP; también se puede revocar permisos específicos)
    \end{block}
\end{frame}

\begin{frame}[fragile]{Ejercicio 2: Configurar PITR en PostgreSQL con Docker}
    \textbf{Enunciado:} Levantar PostgreSQL con archiving WAL y restaurar a un momento antes de un error.
    \begin{block}{Pasos resumidos}
        \begin{lstlisting}[style=bashstyle]
# Crear contenedor con archiving
docker run -d --name pg-pitr -e POSTGRES_PASSWORD=admin -v pgdata:/var/lib/postgresql/data -v walarchive:/wal_archive postgres:15 -c wal_level=replica -c archive_mode=on -c archive_command='cp %p /wal_archive/%f' -c archive_timeout=60

# Simular datos y luego DROP
docker exec -it pg-pitr psql -U postgres -c "CREATE TABLE test (id int); INSERT INTO test VALUES (1); SELECT pg_switch_wal();"

# Restaurar: detener contenedor, preparar recovery, iniciar nuevo contenedor...
        \end{lstlisting}
    \end{block}
\end{frame}

\begin{frame}[fragile]{Ejercicio 3: Integración PostgreSQL + Redis con CDC simulado}
    \textbf{Enunciado:} Usar trigger y worker Python para sincronizar productos con Redis.
    \begin{block}{Solución (ver diapositivas anteriores)}
        (El código ya se mostró en la sección de sincronización).
    \end{block}
\end{frame}

\begin{frame}[fragile]{Ejercicio 4: Despliegue híbrido con Docker Compose}
    \textbf{Enunciado:} Crear docker-compose.yml con PostgreSQL, Redis y una app Python.
    \begin{block}{Solución (ver diapositiva de orquestación)}
        (El YAML ya se mostró).
    \end{block}
\end{frame}

\begin{frame}[fragile]{Ejercicio 5: Detección de data drift con Python}
    \textbf{Enunciado:} Usar test de Kolmogorov-Smirnov para detectar drift.
    \begin{block}{Solución (ver diapositiva de data drift)}
        (El código Python ya se mostró).
    \end{block}
\end{frame}

% ============================================================================
% SECCIÓN 8: GLOSARIO
% ============================================================================
\section{Glosario}

\begin{frame}{Términos Clave}
    \begin{description}
        \item[DRP] Disaster Recovery Plan, plan de recuperación ante desastres.
        \item[PITR] Point-in-Time Recovery, recuperación a un momento exacto.
        \item[WAL] Write-Ahead Log, registro de transacciones en PostgreSQL.
        \item[CDC] Change Data Capture, captura de cambios en bases de datos.
        \item[IaC] Infrastructure as Code, gestión de infraestructura mediante código.
        \item[Snapshot] Imagen instantánea de un sistema de archivos.
        \item[Data Drift] Cambio en la distribución de datos que afecta modelos de IA.
        \item[Egreso] Transferencia de datos fuera de un proveedor cloud (costo).
        \item[Consistencia eventual] Los datos se vuelven consistentes después de un tiempo.
        \item[Políglota] Uso de múltiples tipos de bases de datos en una misma aplicación.
    \end{description}
\end{frame}

% --- Diapositiva final ---
\begin{frame}
    \centering
    \Huge \textbf{¿Preguntas?}
    
    \vspace{1cm}
    \large ¡Gracias por su atención!
\end{frame}

\end{document}